% ===================== 第三阶段：指板、音阶与即兴（第 21-30 课） =====================
\pdfbookmark[1]{第三阶段 · 指板、音阶与即兴}{stg3}
\stagepage{3}{stageC}{指板、音阶与即兴}{第 21--30 课}%
{以 CAGED 指型系统建立指板地图，掌握琶音、人工泛音、五声音阶与 Swing/Shuffle 律动，完成教材 Solo 与基础即兴，并开启《Fade to Black》前奏 Solo 长期作业。}%
{\stgitem{第 21 课 · CAGED 吉他指型系统与指板认识}
 \stgitem{第 22 课 · 指板属性学习}
 \stgitem{第 23 课 · 手型串联}
 \stgitem{第 24 课 · 人工泛音}
 \stgitem{第 25 课 · 教材 Solo 练习}}%
{\stgitem{第 26 课 · 琶音、三和弦与七和弦}
 \stgitem{第 27 课 · 曲目练习与长期作业}
 \stgitem{第 28 课 · 大小调五声音阶手型关联与找调}
 \stgitem{第 29 课 · Swing 与 Shuffle}
 \stgitem{第 30 课 · 五声音阶即兴练习}}
\SetStage{stageC}{第三阶段 · 指板、音阶与即兴}{第 21--30 课}

\pdfbookmark[2]{第 21 课 · CAGED 吉他指型系统与指板认识}{lsn21}
\begin{lessoncard}{21}{CAGED 吉他指型系统与指板认识}
\cardsec{\faBullseye}{教学目标}
\begin{itemize}
\item 理解 CAGED 系统由 C、A、G、E、D 五种开放和弦形状构成。
\item 建立开放和弦形状与封闭把位和弦之间的联系。
\item 能够在不同把位辨认同一和弦的不同 CAGED 指型。
\item 通过根音位置进一步认识指板上的和弦分布。
\end{itemize}
\cardsec{\faIcon{book-open}}{教学内容}
\begin{itemize}
\item C、A、G、E、D 五种和弦形状的基本特征。
\item 开放和弦形状移动到其他把位后的变化。
\item 每个指型中的根音位置及相邻指型之间的衔接。
\item 使用 CAGED 系统在指板上寻找同一和弦的不同位置。
\end{itemize}
\cardsec{\faIcon{pen-fancy}}{课堂练习}
\begin{itemize}
\item 分别辨认五种开放和弦形状及其根音。
\item 给定一个和弦，在指板不同区域寻找对应的 CAGED 指型。
\item 对相邻两个指型进行位置对照与简单切换。
\item 通过根音检查所找到的指型是否正确。
\end{itemize}
\begin{fignote}
\begin{center}
\begin{adjustbox}{max width=\textwidth}
\begin{chordbox}{C}
  \cbmute{6}\cbdot[3]{5}{3}\cbdot[2]{4}{2}\cbopen{3}\cbdot[1]{2}{1}\cbopen{1}
\end{chordbox}\hspace{7mm}%
\begin{chordbox}{A}
  \cbmute{6}\cbopen{5}\cbdot[1]{4}{2}\cbdot[2]{3}{2}\cbdot[3]{2}{2}\cbopen{1}
\end{chordbox}\hspace{7mm}%
\begin{chordbox}{G}
  \cbdot[2]{6}{3}\cbdot[1]{5}{2}\cbopen{4}\cbopen{3}\cbopen{2}\cbdot[3]{1}{3}
\end{chordbox}\hspace{7mm}%
\begin{chordbox}{E}
  \cbopen{6}\cbdot[2]{5}{2}\cbdot[3]{4}{2}\cbdot[1]{3}{1}\cbopen{2}\cbopen{1}
\end{chordbox}\hspace{7mm}%
\begin{chordbox}{D}
  \cbmute{6}\cbmute{5}\cbopen{4}\cbdot[1]{3}{2}\cbdot[3]{2}{3}\cbdot[2]{1}{2}
\end{chordbox}
\end{adjustbox}
\end{center}
\begin{fignotes}
\item 这五个开放形状就是 CAGED 的全部母型：任何封闭把位和弦都是其中某个形状平移后的样子。
\item 学每个形状时先记住根音落在哪根弦——顺着根音就能在指板不同区域找到同一个和弦。
\item 五个形状在指板上首尾衔接：C 型的高音端接着 A 型的低音端，依次循环铺满整条指板。
\end{fignotes}
\end{fignote}
\end{lessoncard}

\pdfbookmark[2]{第 22 课 · 指板属性学习}{lsn22}
\begin{lessoncard}{22}{指板属性学习}
\cardsec{\faBullseye}{教学目标}
\begin{itemize}
\item 理解音符在吉他指板上的重复分布规律。
\item 建立横向移动、纵向换弦与同音异位的基本认识。
\item 能够利用根音、八度和常见位置关系快速定位音符。
\item 为后续手型串联、琶音与音阶学习建立指板地图。
\end{itemize}
\cardsec{\faIcon{book-open}}{教学内容}
\begin{itemize}
\item 六根弦之间的定弦关系与品位移动规律。
\item 同一个音在不同弦、不同品位上的位置。
\item 常见八度指型及根音定位方法。
\item 指板上横向、纵向和区域化观察的区别。
\item 将音名、级数和指型位置联系起来。
\end{itemize}
\cardsec{\faIcon{pen-fancy}}{课堂练习}
\begin{itemize}
\item 给出一个音名，在多个弦组上寻找相同音。
\item 以根音为起点寻找相邻的八度位置。
\item 在限定把位内完成音名或级数定位练习。
\item 使用前一课的 CAGED 根音位置检查指板分布。
\end{itemize}
\begin{fignote}
\begin{center}
\begin{adjustbox}{max width=\textwidth}
\begin{fretboard}[10]
  \draw[stagecur!55, dashed, line width=0.8pt] (4.5,0) -- (6.5,2);
  \draw[stagecur!55, dashed, line width=0.8pt] (4.5,1) -- (6.5,3);
  \draw[stagecur!55, dashed, line width=0.8pt] (4.5,2) -- (7.5,4);
  \draw[stagecur!55, dashed, line width=0.8pt] (4.5,3) -- (7.5,5);
  \fbroot{6}{5}{A}\fbmark{4}{7}{A}
  \fbroot{5}{5}{D}\fbmark{3}{7}{D}
  \fbroot{4}{5}{G}\fbmark{2}{8}{G}
  \fbroot{3}{5}{C}\fbmark{1}{8}{C}
\end{fretboard}
\end{adjustbox}
\end{center}
\begin{fignotes}
\item 常用八度指型：隔一根弦、向高把位挪 2 品就是同名音的高八度（红点到蓝点）。
\item 唯一的例外在 2 弦：从 4 弦或 3 弦出发跨到 2 弦、1 弦时要多挪 1 品（挪 3 品），因为 2 弦定弦低了半音。
\item 记熟「根音＋八度」两点一组，配合 CAGED 根音，整条指板的同名音地图就拼出来了。
\end{fignotes}
\end{fignote}
\end{lessoncard}

\pdfbookmark[2]{第 23 课 · 手型串联}{lsn23}
\begin{lessoncard}{23}{手型串联}
\cardsec{\faBullseye}{教学目标}
\begin{itemize}
\item 将已经学习的不同指型在指板上连接起来。
\item 理解相邻手型之间的重叠音与过渡位置。
\item 减少只会在单一框格内演奏的限制。
\item 提高横向移动、换把和指板整体观察能力。
\end{itemize}
\cardsec{\faIcon{book-open}}{教学内容}
\begin{itemize}
\item 相邻指型的公共音、根音和连接位置。
\item 手型内部演奏与手型之间移动的区别。
\item 使用滑音、换把或保留指完成自然过渡。
\item 将 CAGED、音阶或和弦音位置进行区域连接。
\end{itemize}
\cardsec{\faIcon{pen-fancy}}{课堂练习}
\begin{itemize}
\item 先分别复习两个相邻手型，再标出重叠音。
\item 使用固定上行、下行路线完成两个手型的连接。
\item 在连接点加入滑音或换把练习。
\item 用短乐句往返两个手型，避免机械地只弹完整音阶。
\end{itemize}
\begin{fignote}
\begin{center}
\begin{adjustbox}{max width=\textwidth}
\begin{fretboard}[12]
  % 把位一独有音（第 5 品列）
  \foreach \s in {1,...,6}{\fbmark{\s}{5}{}}
  % 重叠音（蓝点＋金圈）
  \foreach \s/\f in {6/8, 5/7, 4/7, 3/7, 2/8, 1/8}{
    \fbmark{\s}{\f}{}
    \draw[stageC, line width=1.1pt] (\f-0.5,6-\s) circle (0.29);}
  % 把位二独有音（金点）
  \foreach \s/\f in {6/10, 5/10, 4/10, 3/9, 2/10, 1/10}{\fbmark[stageC]{\s}{\f}{}}
  % 3 弦滑音连接箭头
  \draw[stagecur, line width=1.2pt, -{Stealth[length=2.2mm]}]
    (6.85,3.38) .. controls (7.5,3.85) and (8.0,3.85) .. (8.45,3.42);
\end{fretboard}
\end{adjustbox}
\end{center}
\begin{fignotes}
\item 蓝点是第一把位、金点是第二把位；带金圈的音两个把位共用——它们就是过渡的桥。
\item 最顺手的穿越路线：沿 3 弦用滑音从 7 品滑到 9 品（图中箭头），左手顺势整体挪到新把位。
\item 先在各自把位里弹熟，再用三五个音的短乐句往返穿越重叠区，别一上来就跑完整音阶。
\end{fignotes}
\end{fignote}
\end{lessoncard}

\pdfbookmark[2]{第 24 课 · 人工泛音}{lsn24}
\begin{lessoncard}{24}{人工泛音}
\songrow{曲目练习}{乔伊的《Open Fire》}
\cardsec{\faBullseye}{教学目标}
\begin{itemize}
\item 理解人工泛音的基本发声原理。
\item 掌握右手触弦、拨弦与泛音节点之间的配合。
\item 能够在单音及简单乐句中稳定发出人工泛音。
\item 将人工泛音应用到《Open Fire》的相关片段中。
\end{itemize}
\cardsec{\faIcon{book-open}}{教学内容}
\begin{itemize}
\item 实音品位与高八度泛音节点之间的对应关系。
\item 右手触弦手指、拨弦手指与触弦位置的配合。
\item 人工泛音的音量、清晰度和杂音控制。
\item 《Open Fire》相关片段中的人工泛音位置与演奏处理。
\end{itemize}
\cardsec{\faIcon{pen-fancy}}{课堂练习}
\begin{itemize}
\item 在固定音高上寻找并稳定发出人工泛音。
\item 将人工泛音移动到不同琴弦和品位。
\item 对比普通实音与人工泛音的音色和音高关系。
\item 拆分练习《Open Fire》中涉及人工泛音的乐句。
\end{itemize}
\begin{fignote}
\begin{center}
\begin{tikzpicture}[x=1cm,y=1cm]
  \draw[inkGray, line width=1pt] (0,0) -- (7.4,0);
  \draw[line width=2.4pt, inkGray] (0,-0.30) -- (0,0.30);
  \draw[line width=2pt, inkGray!70] (7.4,-0.30) -- (7.4,0.30);
  \draw[stagecur!55, line width=0.8pt]
    (0.9,0) .. controls (1.85,0.62) and (2.8,0.62) .. (3.75,0)
            .. controls (4.7,-0.62) and (5.65,-0.62) .. (6.6,0) -- (7.4,0);
  \draw[stagecur!35, line width=0.8pt, dashed]
    (0.9,0) .. controls (1.85,-0.62) and (2.8,-0.62) .. (3.75,0)
            .. controls (4.7,0.62) and (5.65,0.62) .. (6.6,0) -- (7.4,0);
  \fill[stagecur] (0.9,0) circle (0.13);
  \draw[stagecur, line width=1.1pt] (3.75,0) circle (0.15);
  \draw[stagecur, line width=0.8pt, -{Stealth[length=1.8mm]}] (3.75,1.05) -- (3.75,0.28);
  \node[font=\footnotesize\sffamily, inkGray, anchor=south] at (3.75,1.1) {右手轻触：泛音节点（高 12 品）};
  \node[font=\footnotesize\sffamily, inkGray, anchor=north] at (0.9,-0.45) {左手按实音品};
  \node[font=\footnotesize\sffamily, inkGray, anchor=north] at (6.1,-0.78) {靠近琴桥处拨弦};
  \draw[stagecur, line width=0.8pt, -{Stealth[length=1.8mm]}] (6.1,-0.72) -- (6.35,-0.12);
\end{tikzpicture}
\end{center}
\begin{fignotes}
\item 左手按住实音品位后，弦的振动长度变了，泛音节点跟着移动——节点始终在实音品位高 12 品处。
\item 右手食指轻触节点（对准品丝正上方），拨片同时拨弦，触弦手指随即离开，泛音才能延续。
\item 轻触是关键：只碰弦皮、不压到品丝；触点偏一点，泛音就会发闷或出不来。
\end{fignotes}
\end{fignote}
\end{lessoncard}

\pdfbookmark[2]{第 25 课 · 教材 Solo 练习}{lsn25}
\begin{lessoncard}{25}{教材 Solo 练习}
\songrow{曲目练习}{乔伊的《Open Fire》}
\cardsec{\faBullseye}{教学目标}
\begin{itemize}
\item 通过教材 Solo 综合复习前面学习的演奏技巧。
\item 提高旋律分句、换把、连奏与节奏处理能力。
\item 继续完成《Open Fire》的曲目应用。
\item 建立分段练习、难点循环和完整连接的 Solo 学习方法。
\end{itemize}
\cardsec{\faIcon{book-open}}{教学内容}
\begin{itemize}
\item 教材 Solo 的段落结构、指法与主要技巧。
\item 乐句中的呼吸、重音、延音和连接。
\item 换把、连奏及特殊技巧位置的单独处理。
\item 《Open Fire》相关部分的复习与完整度提升。
\end{itemize}
\cardsec{\faIcon{pen-fancy}}{课堂练习}
\begin{itemize}
\item 将教材 Solo 按乐句拆分，并逐句确定指法。
\item 对换把、连奏和节奏难点进行小范围循环。
\item 先完成准确演奏，再处理乐句力度与表现。
\item 连接教材 Solo 与《Open Fire》中已经学习的内容。
\end{itemize}
\begin{fignote}
\begin{center}
\begin{adjustbox}{max width=\textwidth}
\stepchip{按乐句拆分}\steparrow
\stepchip{逐句定指法}\steparrow
\stepchip{难点小循环}\steparrow
\stepchip{先求准确}\steparrow
\stepchip{再做表情}
\end{adjustbox}
\end{center}
\begin{fignotes}
\item Solo 的最小单位是乐句：像背课文一样一句一句背，指法一旦确定就不再随手换。
\item 「准确」和「表现」分两遍处理：第一遍只管音对、拍对，第二遍再加力度、呼吸和延音。
\item 每句结尾想一下它接谁：句与句的连接处单独过几遍，Solo 才不会一句一停。
\end{fignotes}
\end{fignote}
\end{lessoncard}

\pdfbookmark[2]{第 26 课 · 琶音、三和弦与七和弦}{lsn26}
\begin{lessoncard}{26}{琶音、三和弦与七和弦}
\cardsec{\faBullseye}{教学目标}
\begin{itemize}
\item 理解琶音是和弦音依次演奏形成的旋律线条。
\item 复习三和弦的根音、三音与五音构成。
\item 初步理解七和弦在三和弦基础上加入七音的构成方式。
\item 能够在指板上寻找并演奏基础三和弦与七和弦琶音。
\end{itemize}
\cardsec{\faIcon{book-open}}{教学内容}
\begin{itemize}
\item 三和弦琶音的 1、3、5 构成。
\item 七和弦琶音的 1、3、5、7 构成及基础类型。
\item 和弦手型、和弦音位置与琶音指型之间的联系。
\item 琶音的上行、下行、转位与换弦路线。
\item 在和弦背景上使用琶音突出和弦音。
\end{itemize}
\cardsec{\faIcon{pen-fancy}}{课堂练习}
\begin{itemize}
\item 从已知三和弦中找出根音、三音和五音并依次演奏。
\item 在同一根音上对比三和弦与七和弦琶音。
\item 练习基础琶音的上行、下行和往返连接。
\item 在简单和弦进行中，用对应琶音跟随和弦变化。
\end{itemize}
\begin{fignote}
\begin{center}
\begin{tikzpicture}
  \node[font=\small\sffamily\bfseries, inkGray, left] at (-0.3,0) {三和弦琶音};
  \foreach \x/\n/\c in {1/1/stagecur, 2/3/stagecur!70, 3/5/stagecur!40}{
    \node[fill=\c, text=white, rounded corners=1.2mm, minimum width=0.85cm,
      minimum height=0.6cm, font=\small\sffamily\bfseries] (t\x) at (\x*1.25,0) {\n};}
  \foreach \a/\b in {1/2, 2/3}{
    \draw[inkGray!60, line width=0.7pt, -{Stealth[length=1.6mm]}]
      ([xshift=1.5mm]t\a.east) -- ([xshift=-1.5mm]t\b.west);}
  \node[font=\small\sffamily\bfseries, inkGray, left] at (-0.3,-1.35) {七和弦琶音};
  \foreach \x/\n/\c in {1/1/stagecur, 2/3/stagecur!70, 3/5/stagecur!40, 4/7/stageD!55}{
    \node[fill=\c, text=white, rounded corners=1.2mm, minimum width=0.85cm,
      minimum height=0.6cm, font=\small\sffamily\bfseries] (s\x) at (\x*1.25,-1.35) {\n};}
  \foreach \a/\b in {1/2, 2/3, 3/4}{
    \draw[inkGray!60, line width=0.7pt, -{Stealth[length=1.6mm]}]
      ([xshift=1.5mm]s\a.east) -- ([xshift=-1.5mm]s\b.west);}
\end{tikzpicture}
\end{center}
\begin{fignotes}
\item 琶音就是把和弦音一个一个依次弹出来：不再同时发声，和弦变成了一条旋律线。
\item 三和弦琶音弹 1、3、5；七和弦琶音在此基础上多叠一个七音，色彩立刻更丰富。
\item 上行弹到头就原路下行往返连接，换弦处注意音量均衡——指型直接来自已学的和弦手型。
\end{fignotes}
\end{fignote}
\end{lessoncard}

\pdfbookmark[2]{第 27 课 · 曲目练习与长期作业}{lsn27}
\begin{lessoncard}{27}{曲目练习与长期作业}
\songrow{长期作业}{《Fade to Black》前奏 Solo}
\cardsec{\faBullseye}{教学目标}
\begin{itemize}
\item 综合应用换把、连奏、推弦、揉弦及乐句连接能力。
\item 建立长期曲目练习的拆分、记录与复盘方法。
\item 能够分阶段完成《Fade to Black》前奏 Solo。
\item 通过长期作业提高稳定性、完整度和音乐表达。
\end{itemize}
\cardsec{\faIcon{book-open}}{教学内容}
\begin{itemize}
\item 《Fade to Black》前奏 Solo 的段落划分与技术难点。
\item 指法确定、换把路线和重点技巧位置。
\item 长期作业的练习计划、分段目标与自我检查方法。
\item 从慢速准确到完整连贯的逐步推进方式。
\end{itemize}
\cardsec{\faIcon{pen-fancy}}{课堂练习}
\begin{itemize}
\item 先听辨并划分 Solo 乐句。
\item 按乐句确定指法，对推弦、揉弦和换把单独练习。
\item 选择重点难点进行短循环，不急于从头反复演奏。
\item 记录本周需要完成的具体段落和技术目标。
\end{itemize}
\begin{fignote}
\begin{center}
\begin{adjustbox}{max width=\textwidth}
\stepchip{划分乐句}\steparrow
\stepchip{定指法}\steparrow
\stepchip{重点短循环}\steparrow
\stepchip{记录本周目标}\steparrow
\stepchip{下周复盘推进}
\end{adjustbox}
\end{center}
\begin{fignotes}
\item 长期作业的敌人是「每次都从头弹」：从头弹只巩固前半段，进度会永远卡在中间。
\item 每周只认领一小段：这周的目标写下来（哪几句、到什么速度），下周先复盘再领新段。
\item 推弦、揉弦这类表现技巧单独拎出来练：放进乐句前先在单音上把音准和幅度练稳。
\end{fignotes}
\end{fignote}
\end{lessoncard}

\pdfbookmark[2]{第 28 课 · 大小调五声音阶手型关联与找调}{lsn28}
\begin{lessoncard}{28}{大小调五声音阶手型关联与找调}
\cardsec{\faBullseye}{教学目标}
\begin{itemize}
\item 掌握大小调五声音阶的基本音级构成。
\item 理解关系大小调共用音符与手型的联系。
\item 能够通过根音位置判断当前五声音阶所处的调。
\item 建立不同五声音阶手型之间的关联。
\end{itemize}
\cardsec{\faIcon{book-open}}{教学内容}
\begin{itemize}
\item 大调五声音阶与小调五声音阶的音级构成。
\item 关系大小调的共用音与不同根音中心。
\item 五声音阶手型中的根音位置。
\item 根据伴奏和弦、主音或根音寻找适合的调。
\item 相邻手型之间的重叠音与连接位置。
\end{itemize}
\cardsec{\faIcon{pen-fancy}}{课堂练习}
\begin{itemize}
\item 在同一组音中分别找到大调根音与关系小调根音。
\item 给定调名，在五声音阶手型中标出所有根音。
\item 根据简单伴奏或和弦进行判断调性。
\item 连接两个相邻五声音阶手型，并保持对根音位置的认识。
\end{itemize}
\begin{fignote}
\begin{center}
\begin{adjustbox}{max width=\textwidth}
\begin{fretboard}[10]
  \fbroot{6}{5}{A}\fbmark{6}{8}{C}
  \fbmark{5}{5}{D}\fbmark{5}{7}{E}
  \fbmark{4}{5}{G}\fbroot{4}{7}{A}
  \fbmark{3}{5}{C}\fbmark{3}{7}{D}
  \fbmark{2}{5}{E}\fbmark{2}{8}{G}
  \fbroot{1}{5}{A}\fbmark{1}{8}{C}
\end{fretboard}
\end{adjustbox}
\end{center}
\begin{fignotes}
\item 图为 A 小调五声音阶第一把位（第 5 把位），红色圈点是根音 A——先把三个根音的位置背下来。
\item 同一组音换个中心就是 C 大调五声音阶：落点强调 C 就是大调色彩，强调 A 就是小调色彩。
\item 判断当前调时先听伴奏的根音，再用手型里的根音位置去对——这就是本课「找调」的路径。
\end{fignotes}
\end{fignote}
\end{lessoncard}

\pdfbookmark[2]{第 29 课 · Swing 与 Shuffle}{lsn29}
\begin{lessoncard}{29}{Swing 与 Shuffle}
\cardsec{\faBullseye}{教学目标}
\begin{itemize}
\item 理解直八、Swing 与 Shuffle 之间的区别和联系。
\item 能够听辨平均八分音符与长短八分音符的律动差异。
\item 掌握基础 Shuffle 节奏型，并保持连续稳定的律动。
\item 能够把 Swing 或 Shuffle 感觉应用到基础 Blues 伴奏中。
\end{itemize}
\cardsec{\faIcon{book-open}}{教学内容}
\begin{itemize}
\item 直八原理：每一拍平均分成两个八分音符，两个音的时值基本相等。
\item Swing 原理：把成对八分音符按照三连音框架理解，通常演奏为“前长后短”；教学初期可理解为演奏三连音的第一、第三个位置。实际长短比例会随速度和音乐风格变化，不一定始终固定为严格的二比一。
\item Shuffle 原理：在 Swing 的长短律动基础上，把这种“长、短”关系持续重复，形成清晰而稳定的伴奏型；常见做法同样突出三连音的第一、第三个位置。
\item 两者关系：Swing 更偏向整体的节奏解释和摇摆感觉，Shuffle 更像一种持续、明确、可循环的节奏型或 Groove。
\item 基础十二小节 Blues 进行中的 Shuffle 伴奏应用。
\item 律动练习的重点是稳定、重音和身体感，而不是单纯追求速度。
\end{itemize}
\cardsec{\faIcon{pen-fancy}}{课堂练习}
\begin{itemize}
\item 用口数和拍手对比直八与 Swing 的不同感觉。
\item 在单根空弦上交替演奏直八和 Swing 八分音符。
\item 使用强力和弦或低音弦完成基础 Shuffle 循环。
\item 将 Shuffle 节奏放入简单十二小节 Blues 进行。
\item 配合伴奏连续演奏，重点检查是否越弹越直、长短比例混乱或节奏漂移。
\end{itemize}
\begin{fignote}
\begin{center}
\begin{tikzpicture}[x=1cm,y=1cm]
  \node[font=\small\sffamily\bfseries, inkGray, left] at (-0.25,0.3) {直八};
  \foreach \b in {0,2.2}{
    \draw[inkGray!50, line width=0.7pt] (\b,0) rectangle (\b+2.0,0.6);
    \fill[stagecur!35] (\b+0.03,0.03) rectangle (\b+0.97,0.57);
    \fill[stagecur!35] (\b+1.03,0.03) rectangle (\b+1.97,0.57);}
  \node[font=\scriptsize\sffamily, inkGray!75, right] at (4.5,0.3) {两个八分平均、各占半拍};
  \node[font=\small\sffamily\bfseries, inkGray, left] at (-0.25,-1.0) {Swing};
  \foreach \b in {0,2.2}{
    \draw[inkGray!50, line width=0.7pt] (\b,-1.3) rectangle (\b+2.0,-0.7);
    \draw[inkGray!40, line width=0.5pt, dashed] (\b+0.667,-1.3) -- (\b+0.667,-0.7);
    \draw[inkGray!40, line width=0.5pt, dashed] (\b+1.333,-1.3) -- (\b+1.333,-0.7);
    \fill[stagecur] (\b+0.03,-1.27) rectangle (\b+0.637,-0.73);
    \fill[stagecur!55] (\b+1.363,-1.27) rectangle (\b+1.97,-0.73);}
  \node[font=\scriptsize\sffamily, inkGray!75, right] at (4.5,-1.0) {三连音框架，弹第 1、第 3 位};
\end{tikzpicture}
\end{center}
\begin{fignotes}
\item 直八：一拍平均分成两半，两个音时值基本相等，是流行、摇滚最常见的律动。
\item Swing：把一拍想成三连音（图中虚线三格），只弹第 1、第 3 个位置，自然形成「前长后短」。
\item 长短比例会随速度和风格浮动，不必死抠二比一；把这种律动持续循环下去，就是 Shuffle。
\end{fignotes}
\end{fignote}
\end{lessoncard}

\pdfbookmark[2]{第 30 课 · 五声音阶即兴练习}{lsn30}
\begin{lessoncard}{30}{五声音阶即兴练习}
\cardsec{\faBullseye}{教学目标}
\begin{itemize}
\item 使用已学五声音阶手型完成基础即兴。
\item 能够在伴奏和弦变化时寻找并强调各和弦的根音。
\item 初步建立节奏、停顿、重复与乐句结束意识。
\item 将第 28 课的找调和第 29 课的律动内容应用到即兴中。
\end{itemize}
\cardsec{\faIcon{book-open}}{教学内容}
\begin{itemize}
\item 根据伴奏判断调性并选择对应五声音阶手型。
\item 在每个和弦出现时定位其根音，并将根音作为稳定落点。
\item 使用少量音符进行节奏变化、重复、问答与停顿。
\item 从单一手型即兴逐步过渡到相邻手型连接。
\item 在 Swing 或 Shuffle 伴奏中保持相应律动，而不是只连续跑音阶。
\end{itemize}
\cardsec{\faIcon{pen-fancy}}{课堂练习}
\begin{itemize}
\item 先只使用根音，跟随和弦变化完成落点练习。
\item 限制使用两至三个音，练习不同节奏与停顿。
\item 使用一个五声音阶手型完成短句问答。
\item 在简单伴奏中寻找每个和弦的根音，并把乐句结束在稳定音上。
\item 最后进行短时间自由即兴，检查是否有节奏、呼吸和句子感。
\end{itemize}
\begin{fignote}
\begin{center}
\begin{adjustbox}{max width=\textwidth}
\stepchip{判断调性·选手型}\steparrow
\stepchip{只弹根音找落点}\steparrow
\stepchip{限两三个音造句}\steparrow
\stepchip{加律动与停顿}\steparrow
\stepchip{自由即兴}
\end{adjustbox}
\end{center}
\begin{fignotes}
\item 即兴不是从「随便弹」开始，而是从「只弹根音」开始：根音是每个和弦的安全岛。
\item 音越少句子越像话：限制两三个音逼自己在节奏、重复和停顿上做文章。
\item 乐句要呼吸——弹一句停一拍，问一句答一句；结束音落在稳定音上，句子才算说完。
\end{fignotes}
\end{fignote}
\end{lessoncard}
